\documentclass[11pt,letterpaper]{article}
\usepackage{technical-user-manual}
\setDocTitle{Pipeline Operations, Feedback and Incident Reference}
\setDocSubtitle{Engineering technical reference}
\setDocVersion{1.0}
\setDocStatus{Technical education / proposed reference}
\setDocOwner{CI/CD Engineering}
\setDocDate{30 October 2026}
\setDocSummary{Supplementary engineering practice}
\begin{document}
\makedoctitle
\makedocsummary
\tableofcontents
\clearpage
\section{Purpose, scope and authority}
This is an original modern engineering reference. The 2007 Duvall, Matyas and Glover book informs foundational CI intent but does not prescribe GitHub-specific implementations or 2026 supply-chain standards.

\section{Ownership}
Distinguish platform operator, application maintainer, security reviewer and release approver; record pager/escalation route. Every blocking check has a maintainer and a way to reproduce failures.

Verification: collect the configuration or runtime evidence supporting this control, identify an accountable owner, and record the exception path where the condition cannot presently be met.

\section{Monitoring}
Track pass rate by stage, median and p95 duration, queue time, infrastructure failure, flaky incidence, time to repair a broken mainline, promotion lead time and deployment recovery. Partition by workflow and environment.

Verification: collect the configuration or runtime evidence supporting this control, identify an accountable owner, and record the exception path where the condition cannot presently be met.

\section{Notification design}
Notify change author quickly for PR failures; page the on-call for protected mainline or deployment failures based on severity; report security incidents privately. De-duplicate retry events and record action link and incident ID.

Verification: collect the configuration or runtime evidence supporting this control, identify an accountable owner, and record the exception path where the condition cannot presently be met.

\section{Incident response}
Freeze or revert damaged mainline change, restore service when needed, preserve evidence, identify contributing mechanism, propose preventive test or gate, assign owner and verify the correction on a future run.

Verification: collect the configuration or runtime evidence supporting this control, identify an accountable owner, and record the exception path where the condition cannot presently be met.

\section{Release disposition}
Classify a run as pass, failure, unstable, cancelled or inconclusive. Approved release evidence includes a bound revision and artifact digest; passing historic checks for another revision is insufficient.

Verification: collect the configuration or runtime evidence supporting this control, identify an accountable owner, and record the exception path where the condition cannot presently be met.

\section{Adoption checklist}
\begin{itemize}
\item Identify applicable repositories, application owners and trust boundaries.
\item Define minimum pass criteria and measurable acceptance thresholds.
\item Review workflow permissions, external actions, outputs and retention.
\item Record exceptions with justification, owner and expiry.
\item Revalidate after action, runner or policy changes.
\end{itemize}
\end{document}
